\section{Prompts}\label{app: prompts}

We gather prompts used for \our{} in this section. Our prompt structure is flexible and extensible. To adapt \our{} to a new problem setting, one only needs to define its problem description, function description, and function signature.

\subsection{Common prompts}
The prompt formats are given below. They are used for all COP settings.


% \lstset{
%     backgroundcolor = \color{yellow!10},
%     basicstyle = \ttfamily\small,
%     rulesepcolor= \color{white},
%     breaklines = true,
%     columns = fixed,
%     captionpos = t
% }

% Note: overridden by style
\lstset{
    backgroundcolor=\color{gray!10},  % Light gray background
    basicstyle=\ttfamily\tiny,       % Monospaced font
    frame=single,                     % Adds a frame around the code
    framerule=0.5pt,                  % Frame thickness
    rulecolor=\color{black!30},       % Frame color
    breaklines=true,                  % Automatic line breaking
    breakatwhitespace=true,           % Break lines at whitespaces
    showstringspaces=false,           % Don't show spaces in strings as special character
    columns=flexible,                 % Column spacing
    captionpos=t,                     % Caption position at bottom
    abovecaptionskip=1em,             % Space above caption
    belowcaptionskip=1em,             % Space below caption
    % numbers=left,                     % Line numbers on the left
    numberstyle=\tiny\color{gray},    % Style of the line numbers
}

\renewcommand{\lstlistingname}{Prompt}
\crefname{listing}{Prompt}{Prompts}

\begin{lstlisting}[caption={System prompt for generator LLM.},  label={lst: system prompt for generator LLM}, style=promptstyle]
You are an expert in the domain of optimization heuristics. Your task is to design heuristics that can effectively solve optimization problems.
Your response outputs Python code and nothing else. Format your code as a Python code string: "```python ... ```".    
\end{lstlisting}

\begin{lstlisting}[caption={System prompt for reflector LLM.},  label={lst: system prompt for reflector LLM}, style=promptstyle]
You are an expert in the domain of optimization heuristics. Your task is to give hints to design better heuristics.
\end{lstlisting}


\begin{lstlisting}[caption={Task description.},  label={lst: task description}, style=promptstyle]
Write a {function_name} function for {problem_description}
{function_description}
\end{lstlisting}

\begin{lstlisting}[caption={User prompt for population initialization.},  label={lst: user prompt for population initialization}, style=promptstyle]
{task_description}

{seed_function}

Refer to the format of a trivial design above. Be very creative and give `{func_name}_v2`. Output code only and enclose your code with Python code block: ```python ... ```.

{initial_long-term_reflection}
\end{lstlisting}

\begin{lstlisting}[caption={User prompt for short-term reflection.},  label={lst: user prompt for short-term reflection}, style=promptstyle]
Below are two {function_name} functions for {problem_description}
{function_description}

You are provided with two code versions below, where the second version performs better than the first one.

[Worse code]
{worse_code}

[Better code]
{better_code}

You respond with some hints for designing better heuristics, based on the two code versions and using less than 20 words.
\end{lstlisting}

The user prompt used for short-term reflection in black-box COPs is slightly different from the one used for white-box COPs. We explicitly ask the reflector LLM to infer the problem settings and to give hints about how the node and edge attributes correlate with the black-box objective value.

\begin{lstlisting}[caption={User prompt for short-term reflection on black-box COPs.},  label={lst: user prompt for short-term reflection on black-box COPs}, style=promptstyle]
Below are two {function_name} functions for {problem_description}
{function_description}

You are provided with two code versions below, where the second version performs better than the first one.

[Worse code]
{worse_code}

[Better code]
{better_code}

Please infer the problem settings by comparing two code versions and give hints for designing better heuristics. You may give hints about how edge and node attributes correlate with the black-box objective value. Use less than 50 words.
\end{lstlisting}


\begin{lstlisting}[caption={User prompt for crossover.},  label={lst: user prompt for crossover}, style=promptstyle]
{task_description}

[Worse code]
{function_signature0}
{worse_code}

[Better code]
{function_signature1}
{better_code}

[Reflection]
{short_term_reflection}

[Improved code]
Please write an improved function `{function_name}_v2`, according to the reflection. Output code only and enclose your code with Python code block: ```python ... ```.
\end{lstlisting}

The function signature variables here are used to adjust function names with their versions, which is similar to the design in \cite{romera2023funsearch}. For example, when designing ``heuristics'', the worse code is named ``heuristics\_v0'' while the better code ``heuristics\_v1''. In \cref{lst: user prompt for elitist mutation}, the elitist code is named ``heuristic\_v1''.

\begin{lstlisting}[caption={User prompt for long-term reflection.},  label={lst: user prompt for long-term reflection}, style=promptstyle]
Below is your prior long-term reflection on designing heuristics for {problem_description}
{prior_long-term_reflection}

Below are some newly gained insights.
{new_short-term_reflections}

Write constructive hints for designing better heuristics, based on prior reflections and new insights and using less than 50 words.
\end{lstlisting}

\begin{lstlisting}[caption={User prompt for elitist mutation.},  label={lst: user prompt for elitist mutation}, style=promptstyle]
{task_description}

[Prior reflection]
{long-term_reflection}

[Code]
{function_signature1}
{elitist_code}

[Improved code]
Please write a mutated function `{function_name}_v2`, according to the reflection. Output code only and enclose your code with Python code block: ```python ... ```.
\end{lstlisting}

\subsection{Problem-specific prompt components}

Problem-specific prompt components are given below.
\begin{itemize}
    \item Problem descriptions of all COP settings are given in \cref{tab: problem descriptions}.
    \item The function descriptions of all COP settings are presented in \cref{tab: function descriptions}. The descriptions crafted for black-box settings avoid disclosing any information that could link to the original COP.
    \item The function signatures are gathered in \cref{lst: function signatures used in reevo}.
    \item The seed functions are shown in \cref{lst: seed heuristics used for reevo}. The seed function used for TSP\_constructive is drawn from \cite{liu2023algorithm_evolution}. The seed functions used for black-box ACO settings are expert-designed heuristics \cite{skinderowicz2022_aco_tsp_human_bl, cai2022_aco_cvrp_human_bl, sohrabi2021_aco_op_human_bl, fidanova2020_aco_mkp_human_bl, levine2004_aco_bpp_human_bl}, while those used for while-box ACO settings are trivial all-ones matrices.
    \item The initial long-term reflections for some COP settings are presented in \cref{lst: initial long-term reflections}, while are left empty for the others.
\end{itemize}


    % language=Python,
    % basicstyle=\tiny\ttfamily, 
    % breaklines=true, 
    % emphstyle=\bfseries\color{Rhodamine},
    % commentstyle=\itshape\color{black!50!white},
    % stringstyle=\color{deepgreen},
    % keywordstyle=\ttfamily\color{deepblue},
    % columns=flexible
    
\begin{lstlisting}[caption={Function signatures used in \our{}.},  label={lst: function signatures used in reevo}, style=promptcodestyle, language=Python]
# TSP_NCO
def heuristics(distance_matrix: torch.Tensor) -> torch.Tensor:

# CVRP_NCO
def heuristics(distance_matrix: torch.Tensor, demands: torch.Tensor) -> torch.Tensor:

# DPP_GA_crossover
def crossover(parents: np.ndarray, n_pop: int) -> np.ndarray:

# DPP_GA_mutation
def mutation(population: np.ndarray, probe: int, prohibit: np.ndarray, size: int=100) -> np.ndarray:

# TSP_GLS
def heuristics(distance_matrix: np.ndarray) -> np.ndarray:

# TSP_ACO
def heuristics(distance_matrix: np.ndarray) -> np.ndarray:

# CVRP_ACO
def heuristics(distance_matrix: np.ndarray, coordinates: np.ndarray, demands: np.ndarray, capacity: int) -> np.ndarray:

# OP_ACO
def heuristics(prize: np.ndarray, distance: np.ndarray, maxlen: float) -> np.ndarray:

# MKP_ACO
def heuristics(prize: np.ndarray, weight: np.ndarray) -> np.ndarray:

# BPP_ACO
def heuristics(demand: np.ndarray, capacity: int) -> np.ndarray:

# TSP_ACO (black-box)
def heuristics(edge_attr: np.ndarray) -> np.ndarray:

# CVRP_ACO (black-box)
def heuristics(edge_attr: np.ndarray, node_attr: np.ndarray) -> np.ndarray: # For simplicity, we omit `coordinates` and `capacity` after using capacity to normalize demands, i.e. node_attr

# OP_ACO (black-box)
def heuristics(node_attr: np.ndarray, edge_attr: np.ndarray, node_constraint: float) -> np.ndarray:

# MKP_ACO (black-box)
def heuristics(item_attr1: np.ndarray, item_attr2: np.ndarray) -> np.ndarray:

# BPP_ACO (black-box)
def heuristics(node_attr: np.ndarray, node_constraint: int) -> np.ndarray:

# TSP_constructive
def select_next_node(current_node: int, destination_node: int, unvisited_nodes: set, distance_matrix: np.ndarray) -> int:
\end{lstlisting}


\begin{lstlisting}[
caption={Seed heuristics used for \our{}.},
label={lst: seed heuristics used for reevo}, style=promptcodestyle, language=Python]
# TSP_NCO
def heuristics(distance_matrix: torch.Tensor) -> torch.Tensor:
    distance_matrix[distance_matrix == 0] = 1e5
    K = 100
    # Compute top-k nearest neighbors (smallest distances)
    values, indices = torch.topk(distance_matrix, k=K, largest=False, dim=1)
    heu = -distance_matrix.clone()
    # Create a mask where topk indices are True and others are False
    topk_mask = torch.zeros_like(distance_matrix, dtype=torch.bool)
    topk_mask.scatter_(1, indices, True)
    # Apply -log(d_ij) only to the top-k elements
    heu[topk_mask] = -torch.log(distance_matrix[topk_mask])
    return heu

# CVRP_NCO
def heuristics(distance_matrix: torch.Tensor, demands: torch.Tensor) -> torch.Tensor:
    return torch.zeros_like(distance_matrix)

# DPP_GA_crossover
def crossover(parents: np.ndarray, n_pop: int) -> np.ndarray:
    n_parents, n_decap = parents.shape
    # Split genomes into two halves
    left_halves = parents[:, :n_decap // 2]
    right_halves = parents[:, n_decap // 2:]
    # Create parent pairs
    parents_idx = np.stack([np.random.choice(range(n_parents), 2, replace=False) for _ in range(n_pop)])
    parents_left = left_halves[parents_idx[:, 0]]
    parents_right = right_halves[parents_idx[:, 1]]
    # Create offspring
    offspring = np.concatenate([parents_left, parents_right], axis=1)
    return offspring

# DPP_GA_mutation
def mutation(population: np.ndarray, probe: int, prohibit: np.ndarray, size: int=100) -> np.ndarray:
    n_pop, n_decap = population.shape
    for i in range(n_pop):
        ind = population[i]
        unique_actions = np.unique(population[i])
        if len(unique_actions) < n_decap:
            # Find the indices wherein the action is taken the second time
            dup_idx = []
            action_set = set()
            for j, action in enumerate(ind):
                if action in action_set:
                    dup_idx.append(j)
                action_set.add(action)

            # Mutate the duplicated actions
            infeasible_actions = np.concatenate([prohibit, [probe], unique_actions])
            feasible_actions = np.setdiff1d(np.arange(size), infeasible_actions)
            assert n_decap - len(unique_actions) == len(dup_idx)
            new_actions = np.random.choice(feasible_actions, len(dup_idx), replace=False)
            ind[dup_idx] = new_actions
    return population

# TSP_GLS
def heuristics(distance_matrix: np.ndarray) -> np.ndarray:
    return distance_matrix

# TSP_ACO
def heuristics(distance_matrix: np.ndarray) -> np.ndarray:
    return 1 / distance_matrix

# CVRP_ACO
def heuristics(distance_matrix: np.ndarray, coordinates: np.ndarray, demands: np.ndarray, capacity: int) -> np.ndarray:
    return 1 / distance_matrix

# OP_ACO
def heuristics(prize: np.ndarray, distance: np.ndarray, maxlen: float) -> np.ndarray:
    return prize[np.newaxis, :] / distance

# MKP_ACO
def heuristics(prize: np.ndarray, weight: np.ndarray) -> np.ndarray:
    return prize / np.sum(weight, axis=1)

# BPP_ACO
def heuristics(demand: np.ndarray, capacity: int) -> np.ndarray:
    return np.tile(demand/demand.max(), (demand.shape[0], 1))

# TSP_ACO (black-box)
def heuristics(edge_attr: np.ndarray) -> np.ndarray:
    return np.ones(edge_attr.shape[0])

# CVRP_ACO (black-box)
def heuristics(edge_attr: np.ndarray, node_attr: np.ndarray) -> np.ndarray:
    return np.ones_like(edge_attr)

# OP_ACO (black-box)
def heuristics(node_attr: np.ndarray, edge_attr: np.ndarray, edge_constraint: float) -> np.ndarray:
    return np.ones_like(edge_attr)

# MKP_ACO (black-box)
def heuristics(item_attr1: np.ndarray, item_attr2: np.ndarray) -> np.ndarray:
    n, m = item_attr2.shape
    return np.ones(n,)

# BPP_ACO (black-box)
def heuristics(node_attr: np.ndarray, node_constraint: int) -> np.ndarray:
    n = node_attr.shape[0]
    return np.ones((n, n))

# TSP_constructive
def select_next_node(current_node: int, destination_node: int, unvisited_nodes: set, distance_matrix: np.ndarray) -> int:
    threshold = 0.7
    c1, c2, c3, c4 = 0.4, 0.3, 0.2, 0.1
    scores = {}
    for node in unvisited_nodes:
        all_distances = [distance_matrix[node][i] for i in unvisited_nodes if i != node]
        average_distance_to_unvisited = np.mean(all_distances)
        std_dev_distance_to_unvisited = np.std(all_distances)
        score = c1 * distance_matrix[current_node][node] - c2 * average_distance_to_unvisited + c3 * std_dev_distance_to_unvisited - c4 * distance_matrix[destination_node][node]
        scores[node] = score
    next_node = min(scores, key=scores.get)
    return next_node
\end{lstlisting}


\begin{lstlisting}[caption={Initial long-term reflections},  label={lst: initial long-term reflections}, style=promptstyle]
# White-box COP_ACO
- Try combining various factors to determine how promising it is to select a solution component.
- Try sparsifying the matrix by setting unpromising elements to zero.

# TSP_constructive
- Try look-ahead mechanisms.
\end{lstlisting}



% Create a table of problem descriptions of all problems
\begin{table}[!thp]
    \small
    \centering
    \caption{Problem descriptions used in prompts.}
    \begin{tabular}{l|p{10cm}}
    \toprule
    \multicolumn{1}{c|}{Problem} & \multicolumn{1}{c}{Problem description} \\ \midrule
    TSP\_NCO &  Assisting in solving the Traveling Salesman Problem (TSP) with some prior heuristics. TSP requires finding the shortest path that visits all given nodes and returns to the starting node. \\ \midrule
    CVRP\_NCO & Assisting in solving Capacitated Vehicle Routing Problem (CVRP) with some prior heuristics. CVRP requires finding the shortest path that visits all given nodes and returns to the starting node. Each node has a demand and each vehicle has a capacity. The total demand of the nodes visited by a vehicle cannot exceed the vehicle capacity. When the total demand exceeds the vehicle capacity, the vehicle must return to the starting node. \\ \midrule
    DPP\_GA &  Assisting in solving black-box decap placement problem with genetic algorithm. The problem requires finding the optimal placement of decaps in a given power grid. \\ \midrule
    TSP\_GLS & Solving Traveling Salesman Problem (TSP) via guided local search. TSP requires finding the shortest path that visits all given nodes and returns to the starting node. \\
    \midrule
    TSP\_ACO & Solving Traveling Salesman Problem (TSP) via stochastic solution sampling following "heuristics". TSP requires finding the shortest path that visits all given nodes and returns to the starting node. \\ \midrule
    TSP\_ACO\_black-box & Solving a black-box graph combinatorial optimization problem via stochastic solution sampling following "heuristics". \\
    \midrule
    CVRP\_ACO & Solving Capacitated Vehicle Routing Problem (CVRP) via stochastic solution sampling. CVRP requires finding the shortest path that visits all given nodes and returns to the starting node. Each node has a demand and each vehicle has a capacity. The total demand of the nodes visited by a vehicle cannot exceed the vehicle capacity. When the total demand exceeds the vehicle capacity, the vehicle must return to the starting node. \\ \midrule
    CVRP\_ACO\_black-box & Solving a black-box graph combinatorial optimization problem via stochastic solution sampling following "heuristics".\\
    \midrule
    OP\_ACO & Solving Orienteering Problem (OP) via stochastic solution sampling following "heuristics". OP is an optimization problem where the goal is to find the most rewarding route, starting from a depot, visiting a subset of nodes with associated prizes, and returning to the depot within a specified travel distance. \\ \midrule
    OP\_ACO\_black-box & Solving a black-box graph combinatorial optimization problem via stochastic solution sampling following "heuristics".\\
    \midrule
    MKP\_ACO & Solving Multiple Knapsack Problems (MKP) through stochastic solution sampling based on "heuristics". MKP involves selecting a subset of items to maximize the total prize collected, subject to multi-dimensional maximum weight constraints.\\ \midrule
    MKP\_ACO\_black-box & Solving a black-box combinatorial optimization problem via stochastic solution sampling following "heuristics". \\
    \midrule
    BPP\_ACO & Solving Bin Packing Problem (BPP). BPP requires packing a set of items of various sizes into the smallest number of fixed-sized bins. \\ \midrule
    BPP\_ACO\_black-box & Solving a black-box combinatorial optimization problem via stochastic solution sampling following "heuristics".\\ \midrule
    TSP\_constructive & Solving Traveling Salesman Problem (TSP) with constructive heuristics. TSP requires finding the shortest path that visits all given nodes and returns to the starting node. \\
    \bottomrule
    \end{tabular}
    \label{tab: problem descriptions}
\end{table}


\small{
\begin{longtable}{l|p{10cm}}
    \caption{Function descriptions used in prompts.}\label{tab: function descriptions}\\
    
    \toprule
    \multicolumn{1}{c|}{Problem} & \multicolumn{1}{c}{Function description} \\ 
    \midrule
    \endfirsthead
    
    \multicolumn{2}{c}%
    {{Table \thetable\ continued from previous page}} \\
    \toprule
    \multicolumn{1}{c|}{Problem} & \multicolumn{1}{c}{Function description} \\ 
    \midrule
    \endhead
    
    \midrule
    \multicolumn{2}{r}{{Continued on next page}} \\ \bottomrule
    \endfoot
    
    \bottomrule
    \endlastfoot
    
    TSP\_NCO & The `heuristics` function takes as input a distance matrix and returns prior indicators of how bad it is to include each edge in a solution. The return is of the same shape as the input. The heuristics should contain negative values for undesirable edges and positive values for promising ones. Use efficient vectorized implementations. \\ \midrule
    CVRP\_NCO & The `heuristics` function takes as input a distance matrix (shape: n by n) and a vector of customer demands (shape: n), where the depot node is indexed by 0 and the customer demands are normalized by the total vehicle capacity. It returns prior indicators of how promising it is to include each edge in a solution. The return is of the same shape as the distance matrix. The heuristics should contain negative values for undesirable edges and positive values for promising ones. Use efficient vectorized implementations. \\ \midrule
    DPP\_GA\_crossover &  The `crossover` function takes as input a 2D NumPy array parents and an integer n\_pop. The function performs a genetic crossover operation on parents to generate n\_pop offspring. Use vectorized implementation if possible. \\ \midrule
    DPP\_GA\_mutation & The `mutation` function modifies a given 2D population array to ensure exploration of the genetic algorithm. You may also take into account the feasibility of each individual. An individual is considered feasible if all its elements are unique and none are listed in the prohibited array or match the probe value. Use a vectorized implementation if possible.\newline
    The function takes as input the below arguments:\newline
    - population (np.ndarray): Population of individuals; shape: (P, n\_decap). \newline
    - probe (int): Probe value; each element in the population should not be equal to this value. \newline
    - prohibit (np.ndarray): Prohibit values; each element in the population should not be in this set. \newline
    - size (int): Size of the PDN; each element in the population should be in the range [0, size).  
    \\ \midrule
    TSP\_GLS &  The `heuristics` function takes as input a distance matrix, and returns prior indicators of how bad it is to include each edge in a solution. The return is of the same shape as the input. \\
    \midrule
    TSP\_ACO &  The `heuristics` function takes as input a distance matrix, and returns prior indicators of how promising it is to include each edge in a solution. The return is of the same shape as the input.\\ 
    \midrule
    TSP\_ACO\_black-box &  The `heuristics` function takes as input a matrix of edge attributes with shape `(n\_edges, n\_attributes)`, where `n\_attributes=1` in this case. It computes prior indicators of how promising it is to include each edge in a solution. The return is of the shape of `(n\_edges,)`. \\
    \midrule
    CVRP\_ACO & The `heuristics` function takes as input a distance matrix (shape: n by n), Euclidean coordinates of nodes (shape: n by 2), a vector of customer demands (shape: n), and the integer capacity of vehicle capacity. It returns prior indicators of how promising it is to include each edge in a solution. The return is of the same shape as the distance matrix. The depot node is indexed by 0. \\ 
    \midrule
    CVRP\_ACO\_black-box & The `heuristics` function takes as input a matrix of edge attributes (shape: n by n) and a vector of node attributes (shape: n). A special node is indexed by 0. `heuristics` returns prior indicators of how promising it is to include each edge in a solution. The return is of the same shape as the input matrix of edge attributes.  \\
    \midrule
    OP\_ACO & Suppose `n` represents the number of nodes in the problem, with the depot being the first node. The `heuristics` function takes as input a `prize` array of shape (n,), a `distance` matrix of shape (n,n), and a `max\_len` float which is the constraint to total travel distance, and it returns `heuristics` of shape (n, n), where `heuristics[i][j]` indicates the promise of including the edge from node \#i to node \#j in the solution. \\ 
    \midrule
    OP\_ACO\_black-box & The `heuristics` function takes as input a vector of node attributes (shape: n), a matrix of edge attributes (shape: n by n), and a constraint imposed on the sum of edge attributes. A special node is indexed by 0. `heuristics` returns prior indicators of how promising it is to include each edge in a solution. The return is of the same shape as the input matrix of edge attributes.  \\
    \midrule
    MKP\_ACO & Suppose `n` indicates the scale of the problem, and `m` is the dimension of weights each item has. The constraint of each dimension is fixed to 1. The `heuristics` function takes as input a `prize` of shape (n,), a `weight` of shape (n, m), and returns `heuristics` of shape (n,). `heuristics[i]` indicates how promising it is to include item i in the solution.\\ 
    \midrule
    MKP\_ACO\_black-box &  Suppose `n` indicates the scale of the problem, and `m` is the dimension of some attributes each involved item has. The `heuristics` function takes as input an `item\_attr1` of shape (n,), an `item\_attr2` of shape (n, m), and returns `heuristics` of shape (n,). `heuristics[i]` indicates how promising it is to include item i in the solution.\\
    \midrule
    BPP\_ACO & Suppose `n` represents the number of items in the problem. The heuristics function takes as input a `demand` array of shape (n,) and an integer as the capacity of every bin, and it returns a `heuristics` array of shape (n,n). `heuristics[i][j]` indicates how promising it is to put item i and item j in the same bin. \\ 
    \midrule
    BPP\_ACO\_black-box & Suppose `n` represents the scale of the problem. The heuristics function takes as input an `item\_attr` array of shape (n,) and an integer as a certain constraint imposed on the item attributes. The heuristics function returns a `heuristics` array of shape (n, n). `heuristics[i][j]` indicates how promising it is to group item i and item j.\\
    \midrule
    TSP\_constructive & The select\_next\_node function takes as input the current node, the destination node, a set of unvisited nodes, and a distance matrix, and returns the next node to visit.\\ 
    \bottomrule
    \end{longtable}
}
